\documentclass[10pt,a4paper]{amsart}
\usepackage[margin=19mm]{geometry}
\usepackage{amsmath,amssymb,mathtools}
\usepackage[scr=rsfs,cal=euler]{mathalfa}
\usepackage{xfrac}
\usepackage[unicode,hidelinks,bookmarks=false]{hyperref}
\newcommand{\ie}{i.e.\ }
\newcommand{\cf}{cf.\ }
\newcommand{\resp}{respectively\ }
\newcommand{\Subsection}{\subsection}
\providecommand{\nobreakdash}{\nobreak\hbox{-}}
\setlength{\emergencystretch}{2em}
\multlinegap0pt
\usepackage{amssymb}
\usepackage{float}
\usepackage{tikz}
\usepackage{tcolorbox}
\newtheorem{proposition}{Proposition}
\newcommand{\AAA}{\mathcal{A}}
\newcommand{\Der}{\operatorname{Der}}
\newcommand{\KK}{\mathbb{K}}
\title{A matrix method on the research for the exponents of $2$-multiarrangements}
\author{Shota Maehara}
\date{Source: arXiv:2509.16569v2 (CC BY 4.0)}
\begin{document}
\maketitle

\noindent\emph{Source and licence.} A matrix method on the research for the exponents of $2$-multiarrangements. Version arXiv:2509.16569v2 (CC BY 4.0), distributed by arXiv under the Creative Commons Attribution 4.0 International licence, http://creativecommons.org/licenses/by/4.0/, as recorded in the arXiv licence metadata for this exact version and shown on the abstract page https://arxiv.org/abs/2509.16569v2. Full licence notice: \texttt{licenses/arxiv-2509.16569-CC-BY-4.0.txt}.\par\smallskip

\noindent\textit{Editorial scope.} Counted unit: the article's proposition derivation (source0.tex lines 634--638). The article's complete original proof of it is delivered with it (source0.tex lines 639--649, one \texttt{proof} environment of 11 lines). No mathematics is written, repaired or re-derived: the statement and the proof are the article's own lines, in the article's own notation, and no retained line is substituted or re-flowed.  No further source-local context is needed: every cross-reference of every retained line resolves inside the two delivered spans.  Nothing was omitted from any retained span, so the package carries no omission record.  No retained line contains a citation, so the wrapper carries no bibliography and no bibliography entry.  No retained line contains a figure environment, an image-inclusion command or drawing code, so no image is delivered and the package declares no asset dependency.  Editorial formatting only, in addition: an AMS \texttt{amsart} preamble that restates the article's own declarations and macros used by the retained lines, namely the environments \texttt{proposition} on separate counters, together with the standard packages those lines need.  The retained environments print in this document as Proposition 1 in order of appearance, whereas the article prints the same items with the numbering of its own full document.  The complete native source of the article as distributed in the arXiv e-print for this exact version is bundled unchanged as \texttt{sources/185385/source0.tex}.
        \begin{proposition}\label{proposition:Euler derivation}
    For a derivation $\theta\in\Der S$, 
    $\theta(\alpha_H)\in\alpha_H S$ for infinitely many lines $H\in\KK^2$ 
    if and only if $\theta=f\cdot\theta_E$ by some polynomial $f\in S$. 
    \end{proposition}

\begin{proof}
    If-part is trivial. 
    Then, assume that there exists such a derivation $\theta\neq0$. 
    Since $\theta(\alpha_H)\in\alpha_H S$ for infinitely many lines $H$, 
    we can construct 
    a $2$-dimensional arrangement $\AAA$ such that $|\AAA|-1>\deg(\theta)$ and 
    both $\{\theta_E, \theta\}$ are included in $D(\AAA)$. 
    Since $\exp(\AAA)=(1,|\AAA|-1)$, it follows that 
    $\theta=f\cdot\theta_E$ for some $f\in S$. 
    %which is contradiction. 
\end{proof}

\end{document}
